\documentclass[11pt,a4paper]{article}
\usepackage[utf8]{inputenc}
\usepackage[margin=2cm]{geometry}
\usepackage{xcolor}
\usepackage{enumitem}
\usepackage{tabularx}
\usepackage{booktabs}
\usepackage{fancyhdr}
\usepackage{hyperref}
\usepackage{tcolorbox}
\usepackage{amsmath}
\usepackage{amssymb}
\usepackage{fontawesome5}
\usepackage{listings}
\tcbuselibrary{skins,breakable}

% ── Colors ──
\definecolor{primary}{HTML}{1A237E}
\definecolor{accent}{HTML}{4361EE}
\definecolor{success}{HTML}{2E7D32}
\definecolor{danger}{HTML}{C62828}
\definecolor{warning}{HTML}{E65100}
\definecolor{codebg}{HTML}{F5F5F5}
\definecolor{darkbg}{HTML}{263238}

\hypersetup{colorlinks=true,linkcolor=accent,urlcolor=accent}

% ── Code Listing ──
\lstset{
  backgroundcolor=\color{codebg},
  basicstyle=\ttfamily\small,
  breaklines=true,
  frame=single,
  rulecolor=\color{accent!30},
  keywordstyle=\color{primary}\bfseries,
  commentstyle=\color{success!70},
  stringstyle=\color{danger},
  showstringspaces=false
}

% ── Header/Footer ──
\pagestyle{fancy}
\fancyhf{}
\fancyhead[L]{\small\color{primary}\textbf{Android Security \& Encryption --- Industrial Guide}}
\fancyhead[R]{\small\color{primary}\faIcon{shield-alt}}
\fancyfoot[C]{\thepage}
\renewcommand{\headrulewidth}{1.5pt}
\renewcommand{\headrule}{\hbox to\headwidth{\color{accent}\leaders\hrule height \headrulewidth\hfill}}

% ── Custom Boxes ──
\newtcolorbox{problembox}[1][]{
  enhanced,breakable,colback=danger!6,colframe=danger,
  fonttitle=\bfseries\large\sffamily,title=#1,
  left=8pt,right=8pt,top=6pt,bottom=6pt,
  boxrule=1.5pt,arc=4pt
}
\newtcolorbox{solutionbox}[1][]{
  enhanced,breakable,colback=success!6,colframe=success,
  fonttitle=\bfseries\sffamily,title=#1,
  left=6pt,right=6pt,boxrule=1pt,arc=3pt
}
\newtcolorbox{diagrambox}[1][]{
  enhanced,breakable,colback=codebg,colframe=accent,
  fonttitle=\bfseries\sffamily,title=#1,
  left=6pt,right=6pt,boxrule=1pt,arc=3pt
}
\newtcolorbox{notebox}[1][]{
  enhanced,breakable,colback=warning!8,colframe=warning,
  fonttitle=\bfseries\sffamily,title=#1,
  left=6pt,right=6pt,boxrule=1pt,arc=3pt
}
\newtcolorbox{summarybox}[1][]{
  enhanced,breakable,colback=primary!5,colframe=primary,
  fonttitle=\bfseries\large\sffamily,title=#1,
  left=8pt,right=8pt,top=6pt,bottom=6pt,
  boxrule=1.5pt,arc=4pt
}

\newcommand{\tech}[1]{\textcolor{accent}{\textbf{#1}}}
\newcommand{\danger}[1]{\textcolor{danger}{\textbf{#1}}}
\newcommand{\safe}[1]{\textcolor{success}{\textbf{#1}}}
\newcommand{\stars}[1]{%
  \ifnum#1>0 \faIcon{star}\fi
  \ifnum#1>1 \faIcon{star}\fi
  \ifnum#1>2 \faIcon{star}\fi
  \ifnum#1>3 \faIcon{star}\fi
  \ifnum#1>4 \faIcon{star}\fi
}

\begin{document}

% ══════════════════════════════════════
% TITLE PAGE
% ══════════════════════════════════════
\begin{center}
{\Huge\bfseries\sffamily\color{primary} \faIcon{shield-alt}~Android Security \& Encryption}\\[0.4cm]
{\Large\sffamily\color{accent} Complete Industrial Guide --- Before $|$ During $|$ After Deployment}\\[0.3cm]
{\normalsize\color{gray} 9 Steps: Problem $\rightarrow$ Solution $\rightarrow$ Code $|$ Drone Control App Case Study $|$ Gaurav Bachhav}
\end{center}
\vspace{0.3cm}

\begin{summarybox}[\faIcon{exclamation-triangle}~The Core Problem]
\textbf{Any Android APK can be decompiled} using free tools like JADX, APKTool, or dex2jar.\\
If you hardcode API keys, passwords, or credentials in your code --- \danger{anyone can steal them.}\\[6pt]
\textbf{This guide covers 9 industrial methods} to protect your app, demonstrated on a real drone control application where security is not optional --- it is \textbf{life-critical}.
\end{summarybox}

\tableofcontents
\newpage

% ══════════════════════════════════════════════════
\section{Step 1: Move Secrets to local.properties (BuildConfig Injection)}
% ══════════════════════════════════════════════════

\begin{problembox}[\faIcon{bug}~Problem: Hardcoded Secrets in Source Code]
Secrets written directly in Java code get pushed to GitHub \textbf{and} baked into the APK.
\begin{lstlisting}[language=Java]
// Anyone on GitHub or decompiling the APK sees this:
public static final String API_KEY = "sk-drone-4f8a2b1c-9d3e";
public static final String DB_PASSWORD = "Postgr3s@Pr0d#2026";
\end{lstlisting}
\danger{Risk:} Exposed on GitHub + exposed in APK = double vulnerability.
\end{problembox}

\begin{solutionbox}[\faIcon{check-circle}~Solution: Store in local.properties, inject via Gradle]
Secrets live in \texttt{local.properties} (never pushed to Git). Gradle reads them at build time and injects into \texttt{BuildConfig}.

\textbf{Step A --- Create \texttt{local.properties}:}
\begin{lstlisting}
# This file stays ONLY on your computer
DRONE_API_KEY=sk-drone-4f8a2b1c-9d3e
DB_PASSWORD=Postgr3s@Pr0d#2026
\end{lstlisting}

\textbf{Step B --- Read in \texttt{build.gradle}:}
\begin{lstlisting}[language=Java]
def localProps = new Properties()
localProps.load(new FileInputStream(rootProject.file("local.properties")))

buildConfigField "String", "DRONE_API_KEY",
    "\"${localProps['DRONE_API_KEY']}\""
\end{lstlisting}

\textbf{Step C --- Use in Java:}
\begin{lstlisting}[language=Java]
// No hardcoded value --- reads from BuildConfig
String apiKey = BuildConfig.DRONE_API_KEY;
\end{lstlisting}
\end{solutionbox}

\begin{notebox}[\faIcon{info-circle}~What This Protects \& What It Doesn't]
\begin{tabularx}{\textwidth}{lX}
\toprule
\safe{Protects from} & Secrets appearing on GitHub \\
\danger{Does NOT protect from} & APK decompilation (values are still baked into the APK) \\
\bottomrule
\end{tabularx}
\end{notebox}

% ══════════════════════════════════════════════════
\section{Step 2: R8/ProGuard --- Code Obfuscation}
% ══════════════════════════════════════════════════

\begin{problembox}[\faIcon{bug}~Problem: Decompiled Code is Perfectly Readable]
After decompiling, a hacker sees clean class names, method names, and all debug log messages.
\begin{lstlisting}[language=Java]
// Hacker sees:
public class DroneApiService {
    public void connectToServer() { ... }
    Log.d(TAG, "Password: " + password);  // secret printed!
}
\end{lstlisting}
\end{problembox}

\begin{solutionbox}[\faIcon{check-circle}~Solution: Enable R8 Obfuscation + Strip Logs]
\textbf{Step A --- Enable in \texttt{build.gradle}:}
\begin{lstlisting}[language=Java]
release {
    minifyEnabled true      // Scramble names
    shrinkResources true    // Remove unused code
}
\end{lstlisting}

\textbf{Step B --- Strip logs in \texttt{proguard-rules.pro}:}
\begin{lstlisting}
-assumenosideeffects class android.util.Log {
    public static int v(...);
    public static int d(...);
    public static int i(...);
}
\end{lstlisting}

\textbf{Result --- Hacker now sees:}
\begin{lstlisting}[language=Java]
public class b {
    public void a() { ... }
    // Log lines completely DELETED from release APK
}
\end{lstlisting}
\end{solutionbox}

\begin{diagrambox}[\faIcon{project-diagram}~Before vs After Obfuscation]
\begin{verbatim}
  BEFORE (readable)                    AFTER (scrambled)
  DroneApiService.java                 b.java
    connectToServer()       -->          a()
    sendDroneCommand()      -->          c()
    getTelemetryData()      -->          d()
    Log.d("Password:xyz")  -->          [DELETED]
\end{verbatim}
\end{diagrambox}

\begin{notebox}[\faIcon{info-circle}~Key Fact]
ProGuard only works on \textbf{Java/Kotlin} code. It does \textbf{NOT} protect JSON files, XML files, or assets. That is why Steps 5 and 9 exist.
\end{notebox}

\newpage
% ══════════════════════════════════════════════════
\section{Step 3: EncryptedSharedPreferences --- Secure Local Storage}
% ══════════════════════════════════════════════════

\begin{problembox}[\faIcon{bug}~Problem: User Data Stored as Plain Text on Phone]
Normal \texttt{SharedPreferences} saves data as a readable XML file on the device.
\begin{lstlisting}[language=XML]
<!-- /data/data/com.droneapp/shared_prefs/drone_prefs.xml -->
<string name="password">MyP@ssw0rd123!</string>
<string name="auth_token">eyJhbGciOiJIUzI1...</string>
\end{lstlisting}
\danger{Risk:} Anyone with root access, backup tools, or malware can read this file.
\end{problembox}

\begin{solutionbox}[\faIcon{check-circle}~Solution: Use EncryptedSharedPreferences]
Uses \tech{Android Keystore} (hardware chip) for encryption. Both key names and values are encrypted.

\begin{lstlisting}[language=Java]
// Add dependency: androidx.security:security-crypto:1.1.0-alpha06

String masterKey = MasterKeys.getOrCreate(MasterKeys.AES256_GCM_SPEC);

SharedPreferences securePrefs = EncryptedSharedPreferences.create(
    "secure_drone_prefs",
    masterKey,
    context,
    PrefKeyEncryptionScheme.AES256_SIV,    // encrypt key names
    PrefValueEncryptionScheme.AES256_GCM   // encrypt values
);

// Usage is IDENTICAL to normal SharedPreferences:
securePrefs.edit().putString("password", "MyP@ssw0rd123!").apply();
\end{lstlisting}

\textbf{Stored file now looks like:}
\begin{lstlisting}[language=XML]
<string name="Qp7rT1wYz">xK4vB8cN2jF5hR9mD3...</string>
\end{lstlisting}
\end{solutionbox}

\begin{diagrambox}[\faIcon{project-diagram}~How Encryption Works Internally]
\begin{verbatim}
  Your Code                 Android Keystore (Hardware Chip)
  "save password"  ------>  Master Key encrypts it  ------>  Encrypted XML
  "read password"  <------  Master Key decrypts it  <------  Encrypted XML

  The Master Key NEVER leaves the hardware chip.
  Even if the phone is rooted, the key cannot be extracted.
\end{verbatim}
\end{diagrambox}

% ══════════════════════════════════════════════════
\section{Step 4: Server-Side Proxy --- Remove Secrets from APK Entirely}
% ══════════════════════════════════════════════════

\begin{problembox}[\faIcon{bug}~Problem: API Keys \& DB Credentials Inside the APK]
Even with BuildConfig injection, the actual values are baked into the compiled APK. A skilled hacker can still extract them. Direct database connections from a mobile app expose the \textbf{entire database}.
\begin{lstlisting}[language=Java]
// DANGEROUS: App connects directly to database
String conn = "jdbc:postgresql://db.dronecontrol.com:5432/"
    + "?user=db_admin&password=Postgr3s@Pr0d#2026";
\end{lstlisting}
\end{problembox}

\begin{solutionbox}[\faIcon{check-circle}~Solution: App $\rightarrow$ Your Server $\rightarrow$ Database/APIs]
Move all sensitive credentials to your backend server. App only sends a user authentication token.

\begin{lstlisting}[language=Java]
// App code --- NO secrets, only user token
Request request = new Request.Builder()
    .url("https://your-server.com/api/drone/command")
    .addHeader("Authorization", "Bearer " + userToken)
    .post(body).build();

// Server code (Node.js) --- secrets live HERE
const DRONE_API_KEY = process.env.DRONE_API_KEY;  // environment var
fetch('https://drone-api.com/command', {
    headers: { 'X-API-Key': DRONE_API_KEY }
});
\end{lstlisting}
\end{solutionbox}

\begin{diagrambox}[\faIcon{project-diagram}~Architecture Change]
\begin{verbatim}
  BEFORE (dangerous):
  App ---[API_KEY="secret"]--->  Drone API
  Hacker decompiles --> finds "secret"

  AFTER (safe):
  App ---[user_token]---> YOUR Server ---[API_KEY]--->  Drone API
  Hacker decompiles --> finds NO secret
  API_KEY lives only on YOUR server (hacker can't reach it)
\end{verbatim}
\end{diagrambox}

\begin{notebox}[\faIcon{star}~This is the MOST IMPORTANT step]
\textbf{The most secure secret is one that never exists on the device.} All other steps are layers of defense. This step is the \textbf{only 100\% solution} for API keys and database credentials.
\end{notebox}

\newpage
% ══════════════════════════════════════════════════
\section{Step 5: SSL Pinning --- Secure Network Communication}
% ══════════════════════════════════════════════════

\begin{problembox}[\faIcon{bug}~Problem: Man-in-the-Middle (MITM) Attack]
A hacker on the same WiFi can intercept communication between your app and server using a fake certificate.
\begin{verbatim}
App -----> Hacker (reads everything) -----> Your Server
\end{verbatim}
\end{problembox}

\begin{solutionbox}[\faIcon{check-circle}~Solution: Network Security Config with Certificate Pinning]
Tell your app to \textbf{only trust your specific server's certificate fingerprint}.

\textbf{Create \texttt{res/xml/network\_security\_config.xml}:}
\begin{lstlisting}[language=XML]
<network-security-config>
    <domain-config cleartextTrafficPermitted="false">
        <domain includeSubdomains="true">api.dronecontrol.com</domain>
        <pin-set expiration="2027-06-01">
            <pin digest="SHA-256">YOUR_CERT_HASH_HERE==</pin>
        </pin-set>
    </domain-config>
    <base-config cleartextTrafficPermitted="false" />
</network-security-config>
\end{lstlisting}

\textbf{Link in AndroidManifest.xml:}
\begin{lstlisting}[language=XML]
<application
    android:networkSecurityConfig="@xml/network_security_config" >
\end{lstlisting}

\textbf{Result:} Hacker's fake certificate is rejected. All HTTP traffic is blocked. Only HTTPS to your pinned server is allowed.
\end{solutionbox}

% ══════════════════════════════════════════════════
\section{Step 6: Root \& Tamper Detection --- Runtime Checks}
% ══════════════════════════════════════════════════

\begin{problembox}[\faIcon{bug}~Problem: App Runs on Compromised Devices]
Rooted phones bypass all OS-level security. Modified APKs can contain malicious code. For a drone control app, this is \textbf{physically dangerous}.
\end{problembox}

\begin{solutionbox}[\faIcon{check-circle}~Solution: Check Device Integrity at Startup]
\begin{lstlisting}[language=Java]
// Check 1: Root detection
public static boolean isRooted() {
    String[] paths = {"/system/xbin/su", "/system/bin/su", "/sbin/su"};
    for (String p : paths)
        if (new File(p).exists()) return true;
    return false;
}

// Check 2: Emulator detection
public static boolean isEmulator() {
    return Build.FINGERPRINT.contains("generic")
        || Build.MODEL.contains("Emulator");
}

// Check 3: APK signature verification
public static boolean isTampered(Context ctx) {
    String currentHash = getSignatureHash(ctx);
    return !currentHash.equals("YOUR_KNOWN_HASH");
}
\end{lstlisting}

\textbf{Usage in MainActivity:}
\begin{lstlisting}[language=Java]
if (isRooted() || isEmulator() || isTampered(this)) {
    showAlert("Security Alert! App cannot run on this device.");
    finish();  // Close app
}
\end{lstlisting}
\end{solutionbox}

\begin{diagrambox}[\faIcon{project-diagram}~Startup Security Flow]
\begin{verbatim}
  App Opens --> Rooted? --YES--> Block & Exit
                  |NO
                  v
              Emulator? --YES--> Block & Exit
                  |NO
                  v
              Tampered? --YES--> Block & Exit
                  |NO
                  v
              Debugger? --YES--> Block & Exit
                  |NO
                  v
              ALL CLEAR --> App Runs Normally
\end{verbatim}
\end{diagrambox}

\newpage
% ══════════════════════════════════════════════════
\section{Step 7: Play Integrity API --- Google-Verified Checks}
% ══════════════════════════════════════════════════

\begin{problembox}[\faIcon{bug}~Problem: Local Checks Can Be Bypassed]
Step 6 checks run on the phone --- a skilled hacker can patch the APK to skip them. We need verification from an \textbf{external authority that cannot be fooled}.
\end{problembox}

\begin{solutionbox}[\faIcon{check-circle}~Solution: Google's Play Integrity API]
Google's servers verify the device/app and return a signed token. Verification happens on \textbf{your server via Google} --- impossible to fake.

\begin{lstlisting}[language=Java]
// dependency: com.google.android.play:integrity:1.4.0

IntegrityManager manager = IntegrityManagerFactory.create(context);
manager.requestIntegrityToken(
    IntegrityTokenRequest.builder().setNonce(randomNonce).build()
).addOnSuccessListener(response -> {
    String token = response.token();
    sendTokenToYourServer(token);  // Server verifies with Google
});
\end{lstlisting}
\end{solutionbox}

\begin{diagrambox}[\faIcon{project-diagram}~Play Integrity Flow]
\begin{verbatim}
  App --> Google: "Is this device legit?"
  Google checks --> creates signed TOKEN
  App sends TOKEN --> YOUR Server --> Google API: "Verify this token"
  Google responds: Real device? YES | Genuine app? YES | Rooted? NO
  Server tells App: "Verified. Proceed with drone arming."
\end{verbatim}
\end{diagrambox}

\begin{notebox}[\faIcon{balance-scale}~Step 6 vs Step 7]
\begin{tabularx}{\textwidth}{lXX}
\toprule
& \textbf{Step 6 (Local Checks)} & \textbf{Step 7 (Play Integrity)} \\
\midrule
Who checks? & Your code on the phone & Google's servers \\
Can be bypassed? & Yes, by patching APK & No, verified externally \\
Needs internet? & No & Yes \\
Best for? & Quick first defense & Critical operations (arming drone) \\
\bottomrule
\end{tabularx}
\end{notebox}

% ══════════════════════════════════════════════════
\section{Step 8: NDK Native Code --- Secrets in C Language}
% ══════════════════════════════════════════════════

\begin{problembox}[\faIcon{bug}~Problem: Some Secrets Must Stay on Device]
Google Maps API key must be on the device (the SDK needs it locally). Java strings are trivially extracted from a decompiled APK.
\end{problembox}

\begin{solutionbox}[\faIcon{check-circle}~Solution: Store in C Code (Compiles to Machine Code)]
C/C++ compiles to binary machine code (\texttt{.so} files) --- much harder to reverse-engineer.

\textbf{C file (\texttt{native-keys.c}):}
\begin{lstlisting}[language=C]
#include <jni.h>
JNIEXPORT jstring JNICALL
Java_com_droneapp_security_NativeKeys_getMapsApiKey(
    JNIEnv *env, jobject obj) {
    return (*env)->NewStringUTF(env, "AIzaSyB-YOUR-KEY");
}
\end{lstlisting}

\textbf{Java bridge (\texttt{NativeKeys.java}):}
\begin{lstlisting}[language=Java]
public class NativeKeys {
    static { System.loadLibrary("native-keys"); }
    public static native String getMapsApiKey(); // runs C code
}
\end{lstlisting}
\end{solutionbox}

\begin{diagrambox}[\faIcon{project-diagram}~Java vs C After Decompiling]
\begin{verbatim}
  Java decompiled (readable):
  String KEY = "AIzaSyB-YOUR-KEY";    <-- hacker reads this easily

  C decompiled (machine code):
  0x7f3a: mov  r0, #0x41              <-- meaningless numbers
  0x7f3e: strb r0, [sp, #0x0]         <-- takes hours to decode
  0x7f42: mov  r0, #0x49
\end{verbatim}
\end{diagrambox}

\newpage
% ══════════════════════════════════════════════════
\section{Complete Security Architecture --- Final Overview}
% ══════════════════════════════════════════════════

\begin{summarybox}[\faIcon{trophy}~All 9 Steps at a Glance]
\begin{tabularx}{\textwidth}{clXcc}
\toprule
\textbf{\#} & \textbf{Method} & \textbf{What It Protects} & \textbf{Safety} & \textbf{Difficulty} \\
\midrule
1 & local.properties & Secrets from GitHub & \stars{1} & Easy \\
2 & R8/ProGuard & Code readability + logs & \stars{2} & Easy \\
3 & EncryptedSharedPrefs & User data on phone & \stars{4} & Easy \\
4 & Server-Side Proxy & API keys, DB passwords & \stars{5} & Medium \\
5 & SSL Pinning & Network communication & \stars{4} & Easy \\
6 & Root/Tamper Detection & Compromised devices & \stars{3} & Medium \\
7 & Play Integrity API & Google-verified checks & \stars{4} & Medium \\
8 & NDK Native Code & On-device secrets & \stars{3} & Hard \\
\bottomrule
\end{tabularx}

\vspace{8pt}
\stars{1} = Basic ~|~ \stars{3} = Good ~|~ \stars{5} = Maximum Protection
\end{summarybox}

\begin{diagrambox}[\faIcon{project-diagram}~Complete Security Architecture]
\begin{verbatim}
  BUILD TIME                    DEPLOYMENT                 RUNTIME
  ==========                    ==========                 =======
  local.properties              Play App Signing           Android Keystore
  BuildConfig injection         AAB format                 EncryptedSharedPrefs
  R8/ProGuard obfuscation       Play Integrity API         SSL Pinning
  NDK native secrets            Tamper detection            Server-side proxy
  .gitignore hygiene            Root/emulator checks        Token-based auth

  +-----------+    +----------+    +---------+    +---------+
  | Developer | -> | Build +  | -> | Google  | -> | User's  |
  | writes    |    | Obfuscate|    | Play    |    | Phone   |
  | code      |    | + Sign   |    | Store   |    | (safe!) |
  +-----------+    +----------+    +---------+    +---------+
\end{verbatim}
\end{diagrambox}

\begin{summarybox}[\faIcon{bullseye}~Golden Rules]
\begin{enumerate}[leftmargin=*]
\item \textbf{The best secret is one that never exists on the device} --- use server-side proxy.
\item \textbf{No single method is 100\% safe} --- use multiple layers together.
\item \textbf{ProGuard protects Java only} --- JSON, XML, and assets need separate handling.
\item \textbf{Local checks can be bypassed} --- always verify critical actions via your server.
\item \textbf{Never log secrets} --- use ProGuard to strip all \texttt{Log.d()} in release builds.
\end{enumerate}
\end{summarybox}

\end{document}
